\documentclass[10pt,a4paper]{amsart}
\usepackage[margin=19mm]{geometry}
\usepackage{amsmath,amssymb,mathtools}
\usepackage[scr=rsfs,cal=euler]{mathalfa}
\usepackage{xfrac}
\usepackage[unicode,hidelinks,bookmarks=false]{hyperref}
\newcommand{\ie}{i.e.\ }
\newcommand{\cf}{cf.\ }
\newcommand{\resp}{respectively\ }
\newcommand{\Subsection}{\subsection}
\providecommand{\nobreakdash}{\nobreak\hbox{-}}
\setlength{\emergencystretch}{2em}
\multlinegap0pt
\usepackage{amssymb}
\usepackage{xcolor}
\newtheorem{prop}{Proposition}
\newtheorem{lem}{Lemma}
\title{Towards Bigness equivalence}
\author{F. Laytimi, D. S. Nagaraj}
\date{Source: arXiv:2604.27688v1 (CC BY 4.0)}
\begin{document}
\maketitle

\noindent\emph{Source and licence.} Towards Bigness equivalence. Version arXiv:2604.27688v1 (CC BY 4.0), distributed by arXiv under the Creative Commons Attribution 4.0 International licence, http://creativecommons.org/licenses/by/4.0/, as recorded in the arXiv licence metadata for this exact version and shown on the abstract page https://arxiv.org/abs/2604.27688v1. Full licence notice: \texttt{licenses/arxiv-2604.27688-CC-BY-4.0.txt}.\par\smallskip

\noindent\textit{Editorial scope.} Counted unit: the article's lem lem1 (source0.tex lines 280--282). The article's complete original proof of it is delivered with it (source0.tex lines 284--289, one \texttt{proof} environment of 6 lines). No mathematics is written, repaired or re-derived: the statement and the proof are the article's own lines, in the article's own notation, and no retained line is substituted or re-flowed.  Retained uncounted as the source-local context the proof needs: lines 196--201.  Nothing was omitted from any retained span, so the package carries no omission record.  No retained line contains a citation, so the wrapper carries no bibliography and no bibliography entry.  No retained line contains a figure environment, an image-inclusion command or drawing code, so no image is delivered and the package declares no asset dependency.  Editorial formatting only, in addition: an AMS \texttt{amsart} preamble that restates the article's own declarations and macros used by the retained lines, namely the environments \texttt{prop}, \texttt{lem} on separate counters, together with the standard packages those lines need.  The retained environments print in this document as Proposition 1, Lemma 1 in order of appearance, whereas the article prints the same items with the numbering of its own full document.  The complete native source of the article as distributed in the arXiv e-print for this exact version is bundled unchanged as \texttt{sources/199494/source0.tex}.
\begin{prop}\label{prop2} Let $F$ be a vector bundle and $A$ be an ample line bundle  on $X$ then $F$ is big if and only if 
	$$ {\rm H}^{0}(S^r(F)\otimes A^{-1})\neq (0)$$ 
	for some $r \geq 1,$
	where  $A^{-1}$
	is the dual of the bundle $A.$
\end{prop}

\begin{lem}\label{lem1}
Let $L$ be a big line bundle and $M$ be any line bundle on $X.$ Then for any large enough positve integer $N$ the bundle $L^N \otimes M$ is big. 
\end{lem}

\begin{proof}
 note that $L$ is big implies for any ample line bundle $A$ on $X$  there is an positive integer $n$ and an effective line bundle $E$ on $X$ such that $L^n = E \otimes A.$ 
Now choose a positive  integer $m$ such that $A_1 = A^m\otimes M$ is ample. Hence 
$$X\otimes L^{mn} = E^m (\otimes A^m \otimes M) = E^m \otimes A_1.$$
Since $E^m$ is effective hence from  Proposition(\ref{prop2}) we conclude that $X\otimes L^{mn}$ is big.
\end{proof}

\end{document}
